% Generated draft: requires mathematical and source review.
\documentclass[11pt]{article}
\usepackage[margin=25mm]{geometry}
\usepackage{fontspec}
\setmainfont{Latin Modern Roman}
\usepackage{amsmath,amssymb,mathtools,unicode-math}
\setmathfont{Latin Modern Math}
\usepackage{enumitem,needspace,xurl,hyperref}
\setlength{\parindent}{0pt}
\setlength{\parskip}{5pt}
\setlength{\emergencystretch}{3em}
\newcommand{\N}{\mathbb{N}}
\newcommand{\Z}{\mathbb{Z}}
\newcommand{\Q}{\mathbb{Q}}
\newcommand{\R}{\mathbb{R}}
\newcommand{\C}{\mathbb{C}}
\newcommand{\F}{\mathbb{F}}
\newcommand{\A}{\mathbb{A}}
\newcommand{\PP}{\mathbb{P}}
\newcommand{\E}{\mathbb{E}}
\newcommand{\eps}{\varepsilon}
\newcommand{\dd}{\,\mathrm{d}}
\begin{document}
\section*{UnsolvedMath ID: 20003310}
\textbf{Source problem code:} catalogue entry
\subsection{Definitions and mathematical statement}
The entire assigned statement is ``Question 30 is equivalent to the following question:'' Neither side of the claimed equivalence is included. The aid additionally reports a literal Question/Conjecture 30 label mismatch in the official source. Thus this permanent record has no autonomous conjectural target; the following is the aid's conditional audit of a compact-orbit/minimum-spectrum bridge, not a replacement statement of the missing question.

Let \(d\geq2\), let \(\Lambda\subset\R^d\) be unimodular, and let \(A\) be the positive diagonal determinant-one group. Define
\[
 \operatorname{sys}(\Lambda)=\min_{v\in\Lambda\setminus\{0\}}\|v\|_2,
 \qquad\nu(\Lambda)=\inf_{v\in\Lambda\setminus\{0\}}|v_1\cdots v_d|.
\]
The supplied exact identity is
\[
 \inf_{a\in A}\operatorname{sys}(a\Lambda)
 =\sqrt d\,\nu(\Lambda)^{1/d}.
\]
The infimum may be zero or unattained. The minimum \(\nu\) is constant on the \(A\)-orbit.

In the unimodular lattice space \(X_d=\operatorname{SL}_d(\R)/\operatorname{SL}_d(\Z)\), Mahler's criterion makes
\[
 K_\epsilon=\{\Lambda:\operatorname{sys}(\Lambda)\geq\sqrt d\,\epsilon^{1/d}\}
\]
compact for every \(\epsilon>0\). The identity says that an orbit is contained in \(K_\epsilon\) exactly when its \(\nu\)-value is at least \(\epsilon\).

For any specified family \(\mathscr O\) of compact closed \(A\)-orbit objects, the aid's geometric equivalence can therefore be formulated precisely as
\[
 \forall K\subset X_d\text{ compact},\quad
 \#\{O\in\mathscr O:O\subseteq K\}<\infty
\]
if and only if
\[
 \forall\epsilon>0,\quad
 \#\{O\in\mathscr O:\nu(O)\geq\epsilon\}<\infty.
\]
Here confinement means containment of the entire orbit, not merely meeting a compact set. Compactness of a given \(K\) supplies a uniform positive systole bound for the reverse direction. The missing source statements must be checked before identifying this precise confinement formulation with their wording.

Passing from orbit objects to a numerical positive spectrum \(V\subset\R_{>0}\) requires two additional gates. If \(\nu(\mathscr O)\subseteq V\), finiteness of orbit tails implies finiteness of \(V\cap[\epsilon,\infty)\) only after a coverage hypothesis that every relevant value of \(V\) is realized by an orbit in \(\mathscr O\). Conversely, finite numerical tails imply finite orbit tails only if every value has finitely many orbit objects in its fiber. Neither coverage nor finite fibers follows merely from the compactness of \(K_\epsilon\).

The aid reports abstract countermodels for these omissions: unrepresented numerical values defeat the coverage direction, while infinitely many orbit objects assigned one positive value defeat the finite-fiber direction. These logical countermodels are not claimed to be actual lattice-orbit constructions. The source's alleged numerical reformulation is not reproduced, so its precise spectrum, normalization, and missing arithmetic gates remain unresolved rather than silently supplied.
\subsection{Short English statement}
The assigned record is only an incomplete cross-reference. The supplied audit gives an exact diagonal-lattice identity and a confinement equivalence, then separates the two extra requirements needed to pass to numerical tails: value coverage and finite fibers. It does not verify either missing original statement.
\subsection{Sources}
\begin{itemize}
\item Local UnsolvedMath registry, permanent ID \(20003310\), exact source code \texttt{catalogue entry}; canonical \texttt{aim-other-notes.json}, \texttt{notes[116]}, AIM domain Other, registry category Miscellaneous, source tag \texttt{question}.
\item AIM workshop ``Emerging applications of measure rigidity,'' assigned Question 30 cross-reference: \url{https://aimath.org/WWN/measrigid/measrigid.pdf}. The source text is only the quoted sentence. The aid reports a Question/Conjecture label mismatch; neither referenced statement is supplied. No PDF was fetched.
\item The aid invokes Mahler compactness and an exact diagonal systole/product-minimum identity, but supplies no complete authoritative references or the numerical-spectrum formulation being audited. The original literature field is empty.
\item Research attempt 1: result \texttt{context\_only}, problem status at run \texttt{not\_a\_problem}, rigor label \texttt{proved}, novelty confidence \texttt{low}; classification \texttt{OPEN-TRIAGE}, registry status \texttt{open}, suggested difficulty L2. Created and updated 14 August 2026. These are supplied assessments, not a solution of an absent question.
\item Attribution: Ulam AI / UnsolvedMath Contributors, CC BY 4.0. The machine-generated research aid is not peer reviewed and requires independent verification before citation.
\end{itemize}
\end{document}
